\subsection{Målretting}

\textbf{Zipline bør velge en differensiert målrettingsstrategi der kanal,
budskap og virkemidler tilpasses hvert av de tre prioriterte segmentene. En
udifferensiert tilnærming gir for høy kundeanskaffelseskostnad (CAC), både
fordi det reelle markedet er begrenset til noen få geografiske klynger og
fordi adopsjonsbarrierene varierer betydelig mellom segmentene.}

Målretting er metoden som bringer et budskap fram til det valgte segmentet
\parencite{aspelund2026_forelesning4}. Beslutningen henger sammen med
segmenteringen. Når segmentet må være tilgjengelig for markedsføringstiltak
og kostnaden per oppnådd kunde må stå i forhold til livstidsverdien
\parencite[jf.][]{pfeifer2005}, følger kanalvalget direkte av segmentenes
kjennetegn. Analysen struktureres i tre trinn. Først begrunnes valget mellom
udifferensiert og differensiert tilnærming. Deretter introduseres tre typer
digitale målrettingskanaler (intensjonsbasert, publikumsbasert og
relasjonsbasert), som hver har ulike styrker. Til slutt kobles disse til
Ziplines tre segmenter.

\subsubsection{Udifferensiert vs. differensiert målretting}

Udifferensiert målretting kommuniserer samme budskap til hele markedet og er
kostnadseffektivt for kjennskapsbygging av massemarkedsmerker som Coca-Cola
eller McDonald's \parencite{aspelund2026_forelesning4}. Tilnærmingen har to
svakheter i Ziplines tilfelle. For det første treffer den også forbrukere som
aldri vil kjøpe tjenesten, enten fordi de bor utenfor dronedekningen eller
har for høy adopsjonsbarriere. Disse bomskuddene må regnes inn i CAC for de
kundene som faktisk konverterer, og kostnaden blir raskt uforholdsmessig høy
for en tjeneste hvor det reelle markedet er avgrenset til noen få villastrøk.
For det andre forutsetter udifferensiert kommunikasjon at budskapet fungerer
likt på tvers av segmenter. For Zipline er adopsjonsbarrieren, og dermed hva
som utløser handling, forskjellig for en 35-årig toinntektsfamilie i Bærum og
en 72-åring på Slemdal.

Differensiert målretting lar Zipline bruke ulike budskap og kanaler mot ulike
segmenter, og optimere mot målbare utfall per segment
\parencite{aspelund2026_forelesning4}. Det passer en konkurranseinntreden
hvor segmentene representerer ulike strategiske grupper med distinkte
kjøpskriterier \parencite[jf.][]{oster1999}. Ulempen er høyere kompleksitet,
og at segmentene må være veldefinerte og nåbare gjennom kanaler som tillater
målretting på de relevante variablene. Segmenteringsanalysen har allerede
levert den første forutsetningen. Det neste er å vise at kanalstrukturen
tillater den andre.

\subsubsection{Intensjon, publikum og relasjon}

Digitale markedsføringskanaler kan grupperes i tre typer etter hva
målrettingen bygger på \parencite{aspelund2026_forelesning4}.

\textit{Intensjonsbasert markedsføring} retter seg mot forbrukere i det
øyeblikket de uttrykker en intensjon, typisk gjennom et søk. Google Ads og
lignende plattformer lar annonsøren by på søkeord som signaliserer
kommersiell hensikt. Forbrukeren som søker ``matlevering Nesodden'' har et
udekket behov og uttrykker kjøpsintensjon, noe som gir kanalen høy forventet
konverteringsrate. For Zipline er den særlig relevant der tjenesten dekker
et vakuum i markedet, fordi intensjonen allerede er dannet og friksjonen mot
konvertering er lav.

\textit{Publikumsbasert markedsføring} retter seg mot forhåndsdefinerte
segmenter basert på demografi, geografi, interesser og atferd. Meta og
tilsvarende plattformer tilbyr geografisk målretting ned til postnummer og
interessesignaler som fanger opp livsfase og forbruksmønster. Digital
annonsering oppnår slik en presisjon tradisjonelle medier ikke kan matche,
fordi brukerdata kan knyttes til enkeltindivider \parencite{gallagher1997}.

For Zipline er dette hovedkanalen for kjennskapsbygging. Dronelevering er en
ny kategori, og forbrukere søker sjelden etter tjenester de ikke kjenner til.
Intensjonsbasert målretting har derfor begrenset rekkevidde i
lanseringsfasen, mens publikumsbasert målretting kan aktivere
problemerkjennelse innenfor et geografisk avgrenset dekningsområde. Kanalen
åpner også for at budskapet, ikke bare distribusjonen, tilpasses mottakeren.
\textcite{matz2017} viser at annonser tilpasset psykologiske profiler kan
øke konverteringsraten med opptil 50\,\%. Full psykografisk finjustering
krever et datagrunnlag en ny aktør i Norge ikke har ved lansering, men
demografi- og interessesignaler er tilstrekkelig til å differensiere budskap
per segment.

\textit{Relasjonsbasert markedsføring} bygger på at eksisterende kunder og
nære nettverk rekrutterer nye kunder gjennom anbefaling, lokal synlighet og
lojalitetsprogrammer. I digital form tar dette form av verv-en-venn-%
mekanikk, lokalsamfunnsmarkedsføring og aktivering av opinionsledere. Det
teoretiske grunnlaget er diffusjonsteori, som viser at personlige kilder er
den dominerende informasjonskanalen for tidlige brukere av radikalt nye
produkter, fordi oppfattet risiko reduseres gjennom tillitsbaserte relasjoner
\parencite{rogers2003}. Mekanismen gir kanalen en distinkt kostnadsstruktur.
Den er arbeidsintensiv i etableringsfasen, men marginal CAC faller raskt når
eksisterende kunder selv overtar rekrutteringsarbeidet. For Zipline er den
derfor særlig verdifull i segmenter der personlig tillit og lokal synlighet
veier tyngre enn digital eksponering.

De tre kanaltypene utelukker ikke hverandre. I praksis bygges en portefølje
der hver kanal løser en bestemt oppgave langs kundereisen, og vektingen
varierer mellom segmenter.

\subsubsection{Segmentspesifikk målretting}

\textbf{Segment 1 -- Velstående villaforsteder i Bærum: publikumsbasert
hovedkanal.}
Segmentet er stort nok til å rettferdiggjøre betalt digital
kjennskapsbygging, geografisk veldefinert og digitalt aktivt. Hovedvekten
legges derfor på publikumsbasert markedsføring på Meta-plattformene, med
geografisk avgrensning til postnumrene i Lommedalen, Hosle, Jar og Nesøya,
og interessebasert målretting mot forbrukere som allerede har likt eller
fulgt Foodora, Wolt, Oda eller lignende tjenester. Budskapet i denne kanalen
skal aktivere problemerkjennelse (jf.\ steg 1 i kundereisen), ikke selge
direkte. Intensjonsbasert annonsering gjennom Google Ads fungerer som
supplement på søk knyttet til konkurrenttjenester og takeaway fra
restauranter i Sandvika og Lysaker, hvor Zipline kan forstyrre et
eksisterende kjøpsmønster med en umiddelbar fordel i leveringstid.
Relasjonsbasert aktivitet i dette segmentet tar form av en strukturert
verv-en-venn-mekanikk, fordi synlige naboleveranser i villastrøk er blant de
mest kostnadseffektive bevismidlene selskapet kan utnytte.

\textbf{Segment 2 -- Geografisk isolerte på Nesodden: intensjons- og
relasjonsbasert hovedkanal.}
Segmentet har to kjennetegn som styrer kanalvalget. For det første er
etterspørselen etter matlevering til stede, mens tilbudet ikke er det. Den
primære oppgaven er derfor ikke å skape et behov, men å fange opp et
eksisterende behov. For det andre er lokalsamfunnet lite og kommunikativt
tett, med etablerte digitale og fysiske fellesskap (Nesoddposten, lokale
Facebook-grupper, idrettslag, velforeninger) som gir direkte tilgang til
målgruppen uten mellomledd. De to kjennetegnene peker på hver sin hovedkanal.
Intensjonsbasert markedsføring fanger opp den eksisterende etterspørselen.
Søkeord som ``matlevering Nesodden'' har antagelig lav konkurranse nettopp
fordi ingen andre tilbyr tjenesten, noe som gir lav kostnad per klikk og høy
forventet konverteringsrate. Relasjonsbasert markedsføring utnytter den
tette interne informasjonsflyten gjennom partnerskap med lokale restauranter
og dagligvarebutikker, kombinert med lansering som lokalt nyhetsinnhold.
Publikumsbasert annonsering på Meta spiller en støttende rolle for kjennskap
i perioden rundt lansering, men har lavere marginal nytte fordi den interne
informasjonsflyten i segmentet allerede dekker mye av samme funksjon.

\textbf{Segment 3 -- Eldre velstående i Holmenkollen og Vestre Aker:
relasjonsbasert hovedkanal.}
Dette segmentet har høyest verdipotensial per kunde, men også høyest
adopsjonsbarriere. De store publikumsbaserte plattformene har svakere
dekning blant brukere over 67 år, og intensjonsbasert søkemålretting
forutsetter at forbrukeren selv leter etter tjenesten, noe denne gruppen
sjeldent gjør før de har møtt den andre steder \parencite{ssb2023}.
Målrettingen forskyves derfor mot relasjonsbaserte kanaler. Lokalavis
(Akersposten, Vestkantavisa), samarbeid med lokale butikker, samt
demonstrasjonsarrangementer i Holmenkollen og Slemdal, bruker tillit og
fysisk synlighet som bærende mekanismer. Barn og barnebarn, som ofte er
beslutningspåvirkere for digitale tjenester i denne aldersgruppen, kan nås
gjennom en avgrenset publikumsbasert Meta-kampanje med budskap om å hjelpe
foreldre eller besteforeldre til å prøve tjenesten. Intensjonsbasert
markedsføring inntar en minimal rolle, begrenset til generiske søk om
hjemlevering i Oslo Vest.

\subsubsection{Måling og budsjettprioritering}

Valg av differensiert målretting forutsetter at effekten per segment kan
måles. Hver kanal har etablerte måleparametre (CPC og konverteringsrate for
intensjonsbasert, CPM og klikkrate for publikumsbasert, og vervingsrate for
relasjonsbasert) som gjør det mulig å sammenligne kostnad per oppnådd kunde
mot segmentets estimerte livstidsverdi \parencite[jf.][]{pfeifer2005}.
Pilotfasen i Nesodden, som allerede er forankret i aktivitetsprogrammet,
fungerer dermed også som en strukturert test av kanalmiksen. Observerte
CAC/LTV-forhold fra piloten gir grunnlag for å reallokere budsjett før
Bærum-lanseringen. Målrettingen skal fornyes iterativt når data samler seg,
ikke låses til en initial hypotese.
